Tato sekce shrnuje krátké checklisty pro rychlé nasazení AI, praktické prompt‑šablony a tři stručné ukázky pracovních scénářů vhodné pro pilot.

\subsubsection*{Krátký checklist před spuštěním služby}
\begin{itemize}
  \item Ověřit dostupnost funkcí a nastavit přístupy/role administrátorů \cite{kb_ai_101_tipu_pdf_04e4a115_page_8_1}.
  \item Definovat rozsah činností AI, pravidla lidské validace a eskalace; nastavit pravidla uchovávání a přístupu k datům.
  \item Spustit malý end‑to‑end pilot (generování → odevzdání → AI předopraví → pedagogická revize) a doladit retenci dat \cite{kb_ai_101_tipu_pdf_04e4a115_page_8_1}.
  \item Ověřit licence pro doplňkové zdroje (např. Minecraft Education, Hour of Code) pokud budou použity \cite{kb_ai_101_tipu_pdf_04e4a115_page_15_2}.
\end{itemize}

\subsubsection*{Rychlé šablony promptů (copy‑paste)}
\paragraph{1) Systémový prompt — formativní zpětná vazba}
\begin{verbatim}
SYSTEM: Asistent pro formativní zpětnou vazbu.
Vstupy: student_submission, rubric (3-5 bodů), declared_AI_use.
Výstup: 200-250 slov:
  1) 2 věty povzbuzení
  2) 2 konkrétní kroky ke zlepšení
  3) 1 doporučený milník
Pokud declared_AI_use == "ano", přidej 1 větu co ověřit.
\end{verbatim}

\paragraph{2) Deklarace použití AI (do instrukcí)}
\begin{verbatim}
Deklarace AI:
 - Použil/a jsem AI: Ano/Ne
 - Pokud Ano: nástroj a role (nápověda/editace/generování)
 - Jak ověřeno (max. 50 slov)
 - Sebereflexe (max. 200 slov)
\end{verbatim}

\paragraph{3) Generování adaptivních cvičení}
\begin{verbatim}
Vygeneruj 10 cvičení: lineární rovnice (úroveň: 2. roč. BS).
 - 6 krátkých výpočtů, 4 slovní úlohy.
 - U každého požaduj nahrání postupu.
 - Označ obtížnost (lehký/standardní/obtížný).
\end{verbatim}
(Možnost vyžadovat postup a AI předopravit odpovědi; funkce Vzdělávacích akcelerátorů \cite{kb_ai_101_tipu_pdf_04e4a115_page_8_1}.)

\paragraph{4) Rychlý kvíz (Copilot + Forms)}
\begin{verbatim}
Vytvoř test 10 otázek "Austrálie" pro 8. roč. ZŠ:
 - 4 ABC, 4 pravda/nepravda, 2 otevřené.
 - U každé napiš správnou odpověď a krátké vysvětlení.
\end{verbatim}
(Praktické použití s Microsoft 365 Copilot ve Forms \cite{kb_ai_101_tipu_pdf_04e4a115_page_36_1}.)

\subsubsection*{Krátké ukázky pracovních scénářů}
\paragraph{A) Adaptivní matematická praxe}
Učitel nastaví generátor (téma, obtížnost, požadavek na postup), studenti odevzdají postup; AI předběžně označí chyby a shrne časté potíže, učitel provede finální kontrolu \cite{kb_ai_101_tipu_pdf_04e4a115_page_8_1}.

\paragraph{B) Rychlá in‑class písemka}
Učitel vygeneruje a upraví písemku, nasdílí v LMS; Copilot přidá vysvětlení a zrychlí přípravu \cite{kb_ai_101_tipu_pdf_04e4a115_page_36_1}.

\paragraph{C) Seznámení s principy AI}
Krátké workshopy využijí Minecraft Education a Hour of Code pro demonstraci principů ML a interaktivní seznámení studentů \cite{kb_ai_101_tipu_pdf_04e4a115_page_15_2}.

\subsubsection*{Závěrečné doporučení}
Přizpůsobte šablony místnímu kurikulu a vždy validujte výstupy AI lidskou kontrolou.